% !TEX TS-program = xelatex
% !TEX encoding = UTF-8 Unicode
% ================================================
%  工作周报
%  编译: xelatex → bibtex → xelatex → xelatex
% ================================================

\documentclass{sjtuweekly}
\usepackage{booktabs}
\usepackage{tabularx}
\usepackage{array}

% ========== 图表编号: 独立递增 (图1, 表1, ...) ==========
\renewcommand{\thefigure}{\arabic{figure}}
\renewcommand{\thetable}{\arabic{table}}

% ========== 参考文献库 ==========
\bibfile{bib/database}

% ========== 基本信息 ==========
\member{霍锦龙}
% \weeknum{xx}

\begin{document}

\makeweeklytitle   % 封面（第1页）

% ================================================
%  正文从第2页开始
% ================================================

\section{一、本周工作}

% \subsection*{1.\quad MoE下Pretrigger影响讨论}

\textbf{真实模型权重 + OCS 仿真结合。}
此前分别完成了两件事：(1) 基于合成专家的 MoE+OCS 通信测试床；(2) Qwen3.6-35B-A3B-4bit router gate 的 LoRA 微调与路由 trace 捕获。本周将二者合并，进行 OCS 仿真，替换合成专家，在真实路由 trace 驱动并讨论 8 专家; 32 专家的完整 OCS配置; 双批次重叠，重配置隐藏于计算窗口. 主要目的是在真实模型权重和路由分布下评估 OCS 预触发的通信收益与电路复用率，解决此前合成专家的结论可信度问题。合成专家完成 OCS与 EPS 对比：

\begin{table}[htbp]
\centering
\small
\begin{tabular}{lccc}
\toprule
\textbf{指标} & \textbf{EPS} & \textbf{OCS} & \textbf{变化} \\
\midrule
通信时间        & 23,269 µs  & 17,091 µs  & --26.6\% \\
OCS 额外开销    & —          & 333 µs     & — \\
电路复用率      & —          & 96.7\%     & — \\
重配置时间      & —          & 150 µs (3 次冷启动) & — \\
\bottomrule
\end{tabular}
\caption{合成专家 + OCS 对比：电路建立后 87/90 走光路，较EPS 延迟优势明显。}
\label{tab:ocs_eps}
\end{table}

结论：OCS 电路复用率高；OCS pipeline 模式下重配置与计算重叠，DBO 模式下完全隐藏。

% ================================================
\section{二、工作计划}

\begin{enumerate}
  \item 在 Qwen+OCS 测试床上完成多 batch、多拓扑的完整 benchmark，重标定真实权重下 OCS 收益.
  \item 对比 LoRA 微调前后的路由 trace 在 OCS 仿真中的表现差异，验证路由可学习性对 OCS 预触发的实际增益.
  \item 结合 vLLM prefill/decode 阶段通信模式，讨论 OCS 预触发在推理服务场景中的适配粒度.
\end{enumerate}

% ================================================
\printbib

\end{document}
